% Question bank of the linear algebra practice page:
%   https://prasadsachchidanand.github.io/miscellaneous/linear-algebra-practice/
%
% The page reads this file as it is, so to add or change a question, edit this file and push. Before
% pushing you can check it with
%   node check-questions.mjs
% (run in this folder): it lists any question the page could not read. Compiling the file with pdflatex
% also works, and gives a PDF of the questions with answers and solutions.
%
% ---------------------------------------------------------------------------------------------- FORMAT
%
% \topic{id}{name}{weeks}         starts a topic; the questions below it belong to it
%
% \begin{question}[marks=2]       the [...] is optional; see the keys below
%   The question, with maths in $...$ or \[...\].
%   \begin{options}               multiple choice: 2 to 4 options
%     \option  a wrong option
%     \correct the right option   (several \correct make it a multiple-select question)
%   \end{options}
%   \begin{solution}
%     The worked solution. It can refer to the options as (A), (B), (C), (D); they are never shuffled.
%   \end{solution}
% \end{question}
%
% For a numerical-answer question, write \answer{...} in place of the options:
%   \answer{2}               the answer is 2 (decimals such as -0.75 are fine; not fractions)
%   \answer{6.70 to 6.72}    any number from 6.70 to 6.72 is accepted
%
% Keys in [...], all optional:
%   marks=2       1 is the default
%   type=msq      only needed for a multiple-select question with just one \correct option;
%                 otherwise the type follows from the question (\answer: numerical; one \correct:
%                 MCQ; more than one \correct: MSQ)
%   id=sys-9      the page remembers each visitor's results by this id. A question without one gets
%                 an id made from its text, so editing that text later just restarts its record.
%
% On the page, maths is set by MathJax (amsmath and amssymb commands are fine). In the text, \emph,
% \textbf, \texttt, \href, \url, itemize and enumerate lists, --, ---, ~, ``quotes'' and blank lines
% between paragraphs are converted; the checker warns about any other command.
% ------------------------------------------------------------------------------------------------------

\documentclass[11pt]{article}
\usepackage{practice-questions}   % practice-questions.sty, in this folder
% \showanswersfalse               % uncomment for a PDF without answers and solutions

\begin{document}
{\Large\bfseries Linear algebra practice questions}

% ==============================================================================================
\topic{sys}{Linear systems and row reduction}{Weeks 1--2}

\begin{question}[id=sys-1]
The system $x + y + z = 1,\; x + 2y + 3z = 2,\; x + 3y + 5z = k$ is consistent if and only if
\begin{options}
  \option $k = 1$
  \option $k = 2$
  \correct $k = 3$
  \option it is consistent for every $k$
\end{options}
\begin{solution}
Subtract the first equation from the others: $y + 2z = 1$ and $2y + 4z = k - 1$. Twice the first of these gives $2y + 4z = 2$, so the system is consistent exactly when $k - 1 = 2$, that is, $k = 3$.
\end{solution}
\end{question}

\begin{question}[id=sys-2]
For which value of $a$ does the system $x + 2y = 3,\; 2x + ay = 6$ have infinitely many solutions?
\begin{options}
  \option $a = 2$
  \correct $a = 4$
  \option $a = 6$
  \option for no value of $a$
\end{options}
\begin{solution}
$R_2 \to R_2 - 2R_1$ gives $(a - 4)y = 0$. If $a = 4$ this is $0 = 0$, so $y$ is free and there are infinitely many solutions; if $a \ne 4$, then $y = 0$, $x = 3$ is the only solution.
\end{solution}
\end{question}

\begin{question}[id=sys-3]
How many free variables does the system $Ax = 0$ have, where
\[
  A = \begin{bmatrix} 1 & 2 & 0 & 3 \\ 2 & 4 & 1 & 7 \\ 3 & 6 & 1 & 10 \end{bmatrix}?
\]
\answer{2}
\begin{solution}
Row reduction gives $\begin{bmatrix} 1 & 2 & 0 & 3 \\ 0 & 0 & 1 & 1 \\ 0 & 0 & 0 & 0 \end{bmatrix}$: pivots in columns 1 and 3. The variables $x_2, x_4$ are free, so there are $4 - 2 = 2$ free variables.
\end{solution}
\end{question}

\begin{question}[id=sys-4, type=msq, marks=2]
Let $A$ be a $3 \times 5$ real matrix. Which of the following are always true?
\begin{options}
  \correct $Ax = 0$ has a nontrivial solution.
  \option $Ax = b$ is consistent for every $b \in \mathbb{R}^3$.
  \correct The columns of $A$ are linearly dependent.
  \correct $A$ has at most 3 pivot positions.
\end{options}
\begin{solution}
There is at most one pivot in each row, so at most 3 pivots (D). Then at least $5 - 3 = 2$ of the 5 variables are free, so $Ax = 0$ has nontrivial solutions (A), which means the columns are dependent (C). (B) fails, for instance, for the zero matrix.
\end{solution}
\end{question}

\begin{question}[id=sys-5]
The reduced row echelon form of $\begin{bmatrix} 1 & 2 \\ 3 & 4 \end{bmatrix}$ is
\begin{options}
  \correct $\begin{bmatrix} 1 & 0 \\ 0 & 1 \end{bmatrix}$
  \option $\begin{bmatrix} 1 & 2 \\ 0 & 1 \end{bmatrix}$
  \option $\begin{bmatrix} 1 & 2 \\ 0 & -2 \end{bmatrix}$
  \option $\begin{bmatrix} 1 & 0 \\ 0 & 0 \end{bmatrix}$
\end{options}
\begin{solution}
The determinant is $4 - 6 = -2 \ne 0$, so the matrix is invertible and its reduced echelon form is $I$. (Option C is an echelon form, but not the reduced one.)
\end{solution}
\end{question}

\begin{question}[id=sys-6, marks=2]
A system of 4 linear equations in 6 unknowns
\begin{options}
  \option always has a solution
  \correct has infinitely many solutions whenever it has a solution
  \option has exactly one solution whenever it has a solution
  \option never has a solution
\end{options}
\begin{solution}
At most 4 of the 6 variables can be pivot variables, so a consistent system has at least 2 free variables and hence infinitely many solutions. It need not be consistent: a system that contains the equations $x_1 = 0$ and $x_1 = 1$ has no solution.
\end{solution}
\end{question}

\begin{question}[id=sys-7, marks=2]
The system with augmented matrix
\[
  \left[\begin{array}{ccc|c} 1 & -1 & 2 & 3 \\ 2 & 1 & 1 & 3 \\ 1 & 2 & -1 & k \end{array}\right]
\]
is consistent. Then $k =$
\answer{0}
\begin{solution}
$R_2 \to R_2 - 2R_1$ and $R_3 \to R_3 - R_1$ give the rows $[0\;\, 3\; -3 \mid -3]$ and $[0\;\, 3\; -3 \mid k - 3]$. Then $R_3 \to R_3 - R_2$ gives $[0\;\, 0\;\, 0 \mid k]$, so the system is consistent exactly when $k = 0$.
\end{solution}
\end{question}

\begin{question}[id=sys-8]
Which of the following matrices is in reduced row echelon form?
\begin{options}
  \correct $\begin{bmatrix} 1 & 0 & 2 \\ 0 & 1 & 3 \\ 0 & 0 & 0 \end{bmatrix}$
  \option $\begin{bmatrix} 1 & 2 & 0 \\ 0 & 0 & 1 \\ 0 & 1 & 0 \end{bmatrix}$
  \option $\begin{bmatrix} 0 & 1 & 0 \\ 1 & 0 & 0 \\ 0 & 0 & 1 \end{bmatrix}$
  \option $\begin{bmatrix} 1 & 0 & 0 \\ 0 & 2 & 0 \\ 0 & 0 & 1 \end{bmatrix}$
\end{options}
\begin{solution}
In (B) and (C) a leading entry is to the left of the one in the row above. In (D) a pivot is 2, not 1. Matrix (A) has leading 1s, each with zeros above and below, moving to the right, and its zero row at the bottom.
\end{solution}
\end{question}

\begin{question}[id=sys-9, marks=2]
The system $x + y + z = 3,\; x + 2y + 3z = 6,\; x + 4y + \lambda z = \mu$ has infinitely many solutions when
\begin{options}
  \correct $\lambda = 7,\ \mu = 12$
  \option $\lambda = 7,\ \mu \ne 12$
  \option $\lambda \ne 7$
  \option $\lambda = 12,\ \mu = 7$
\end{options}
\begin{solution}
Subtracting the first equation from the others gives $y + 2z = 3$ and $3y + (\lambda - 1)z = \mu - 3$; subtracting 3 times the first of these from the second gives $(\lambda - 7)z = \mu - 12$. If $\lambda \ne 7$, the solution is unique; if $\lambda = 7$ and $\mu \ne 12$, there is none; if $\lambda = 7$ and $\mu = 12$, $z$ is free and there are infinitely many.
\end{solution}
\end{question}

% ==============================================================================================
\topic{mat}{Matrix algebra and inverses}{Week 3}

\begin{question}[id=mat-1]
If $A$ and $B$ are invertible $n \times n$ matrices, then $(AB)^{-1}$ equals
\begin{options}
  \option $A^{-1}B^{-1}$
  \correct $B^{-1}A^{-1}$
  \option $AB^{-1}$
  \option $BA^{-1}$
\end{options}
\begin{solution}
$(AB)(B^{-1}A^{-1}) = A(BB^{-1})A^{-1} = AA^{-1} = I$. The order reverses, as for putting on socks and then shoes.
\end{solution}
\end{question}

\begin{question}[id=mat-2]
The $(1,2)$ entry of the inverse of $\begin{bmatrix} 2 & 1 \\ 5 & 3 \end{bmatrix}$ is
\answer{-1}
\begin{solution}
The determinant is $6 - 5 = 1$, so the inverse is $\begin{bmatrix} 3 & -1 \\ -5 & 2 \end{bmatrix}$ (swap the diagonal entries, negate the others, divide by the determinant). Its $(1,2)$ entry is $-1$.
\end{solution}
\end{question}

\begin{question}[id=mat-3, type=msq, marks=2]
For $n \times n$ matrices $A$ and $B$ ($n \ge 2$), which of the following are always true?
\begin{options}
  \option $(A + B)^2 = A^2 + 2AB + B^2$
  \correct $(AB)^T = B^T A^T$
  \option If $AB = 0$, then $A = 0$ or $B = 0$.
  \correct $\operatorname{tr}(AB) = \operatorname{tr}(BA)$
\end{options}
\begin{solution}
(A) needs $AB = BA$: in general $(A+B)^2 = A^2 + AB + BA + B^2$. (C) fails for $A = B = \begin{bmatrix} 0 & 1 \\ 0 & 0 \end{bmatrix}$. (B) and (D) always hold: $\operatorname{tr}(AB) = \sum_{i,j} a_{ij}b_{ji} = \operatorname{tr}(BA)$.
\end{solution}
\end{question}

\begin{question}[id=mat-4, marks=2]
Let $A$ be an $n \times n$ matrix with $A^2 = A$ and $A \ne I$. Then
\begin{options}
  \option $A$ is invertible
  \correct $A$ is not invertible
  \option $A = 0$
  \option $A^2 = I$
\end{options}
\begin{solution}
If $A$ were invertible, multiplying $A^2 = A$ by $A^{-1}$ would give $A = I$. So $A$ is not invertible. It need not be 0: $\begin{bmatrix} 1 & 0 \\ 0 & 0 \end{bmatrix}$ satisfies $A^2 = A$.
\end{solution}
\end{question}

\begin{question}[id=mat-5]
If $A$ is $3 \times 4$ and $B$ is $4 \times 2$, the size of $(AB)^T$ is
\begin{options}
  \option $3 \times 2$
  \correct $2 \times 3$
  \option $4 \times 4$
  \option it is not defined
\end{options}
\begin{solution}
$AB$ is $3 \times 2$, so its transpose is $2 \times 3$.
\end{solution}
\end{question}

\begin{question}[id=mat-6, marks=2]
Let $A = \begin{bmatrix} 1 & 2 & 0 \\ 0 & 1 & 3 \\ 0 & 0 & 1 \end{bmatrix}$. The sum of all the entries of $A^{-1}$ is
\answer{4}
\begin{solution}
Write $A = I + N$ with $N = \begin{bmatrix} 0 & 2 & 0 \\ 0 & 0 & 3 \\ 0 & 0 & 0 \end{bmatrix}$. Then $N^3 = 0$, so $A^{-1} = I - N + N^2 = \begin{bmatrix} 1 & -2 & 6 \\ 0 & 1 & -3 \\ 0 & 0 & 1 \end{bmatrix}$, whose entries add up to $4$. (Or row reduce $[A \mid I]$.)
\end{solution}
\end{question}

\begin{question}[id=mat-7]
For an $n \times n$ matrix $A$, which of the following is \emph{not} equivalent to "$A$ is invertible"?
\begin{options}
  \option $Ax = 0$ has only the trivial solution.
  \option The columns of $A$ span $\mathbb{R}^n$.
  \option $A$ has $n$ pivot positions.
  \correct $Ax = b$ has a solution for at least one $b \in \mathbb{R}^n$.
\end{options}
\begin{solution}
(D) holds for every matrix: take $b = 0$. (A), (B) and (C) are all in the Invertible Matrix Theorem.
\end{solution}
\end{question}

\begin{question}[id=mat-8, marks=2]
Let $A$ be a nonzero $2 \times 2$ matrix with $A^2 = 0$. Then
\begin{options}
  \option $\operatorname{rank} A = 0$
  \correct $\operatorname{rank} A = 1$
  \option $\operatorname{rank} A = 2$
  \option $A$ is invertible
\end{options}
\begin{solution}
$A \ne 0$ gives rank at least 1. If $A$ were invertible, $A^2 = 0$ would give $A = 0$; so rank $A \le 1$. Hence rank $A = 1$, as for $\begin{bmatrix} 0 & 1 \\ 0 & 0 \end{bmatrix}$.
\end{solution}
\end{question}

\begin{question}[id=mat-9]
Let $A = \begin{bmatrix} 1 & 1 \\ 0 & 1 \end{bmatrix}$. The $(1,2)$ entry of $A^{10}$ is
\answer{10}
\begin{solution}
By induction, $A^n = \begin{bmatrix} 1 & n \\ 0 & 1 \end{bmatrix}$: if it holds for $n$, then $A^{n+1} = A^nA = \begin{bmatrix} 1 & n + 1 \\ 0 & 1 \end{bmatrix}$. So the entry is 10.
\end{solution}
\end{question}

% ==============================================================================================
\topic{vs}{Vector spaces, subspaces and span}{Week 4}

\begin{question}[id=vs-1, type=msq, marks=2]
Which of the following are subspaces of $\mathbb{R}^3$?
\begin{options}
  \correct $\{(x, y, z) : x + y + z = 0\}$
  \option $\{(x, y, z) : x + y + z = 1\}$
  \option $\{(x, y, z) : xy = 0\}$
  \correct $\{(x, y, z) : x = 2y = 3z\}$
\end{options}
\begin{solution}
(A) is a plane through 0 and (D) is a line through 0 (two linear equations), so both are subspaces. (B) does not contain 0. (C) is not closed under addition: $(1,0,0)$ and $(0,1,0)$ are in it, but their sum $(1,1,0)$ is not.
\end{solution}
\end{question}

\begin{question}[id=vs-2]
Which of the following is a subspace of the space $M_2(\mathbb{R})$ of real $2 \times 2$ matrices?
\begin{options}
  \option the invertible matrices
  \option the matrices with determinant 0
  \correct the symmetric matrices
  \option the matrices with trace 1
\end{options}
\begin{solution}
Sums and scalar multiples of symmetric matrices are symmetric, and 0 is symmetric. The invertible matrices do not contain 0; the matrices with trace 1 do not contain 0; and $\begin{bmatrix} 1 & 0 \\ 0 & 0 \end{bmatrix} + \begin{bmatrix} 0 & 0 \\ 0 & 1 \end{bmatrix} = I$ shows that the singular matrices are not closed under addition.
\end{solution}
\end{question}

\begin{question}[id=vs-3]
The span of $\{(1, 2, 3), (2, 4, 6)\}$ in $\mathbb{R}^3$ is
\begin{options}
  \option $\{0\}$
  \correct a line through the origin
  \option a plane through the origin
  \option all of $\mathbb{R}^3$
\end{options}
\begin{solution}
$(2, 4, 6) = 2(1, 2, 3)$, so the span is the set of multiples of $(1, 2, 3)$: a line through the origin.
\end{solution}
\end{question}

\begin{question}[id=vs-4]
Let $\mathbb{P}_3$ be the space of real polynomials of degree at most 3. The dimension of $\{p \in \mathbb{P}_3 : p(1) = 0\}$ is
\answer{3}
\begin{solution}
$\dim \mathbb{P}_3 = 4$ and $p \mapsto p(1)$ is a nonzero linear map onto $\mathbb{R}$, so its kernel has dimension $4 - 1 = 3$. A basis: $x - 1, (x - 1)^2, (x - 1)^3$.
\end{solution}
\end{question}

\begin{question}[id=vs-5, marks=2]
For which value of $h$ is $(1, h, 5)$ in the span of $(1, 2, 3)$ and $(0, 1, 1)$?
\begin{options}
  \option $h = 2$
  \option $h = 3$
  \correct $h = 4$
  \option $h = 5$
\end{options}
\begin{solution}
We need $a(1, 2, 3) + b(0, 1, 1) = (a, 2a + b, 3a + b) = (1, h, 5)$. So $a = 1$, then $3 + b = 5$ gives $b = 2$, and $h = 2a + b = 4$.
\end{solution}
\end{question}

\begin{question}[id=vs-6, type=msq]
Which of the following subsets of the vector space of all functions $f : \mathbb{R} \to \mathbb{R}$ are subspaces?
\begin{options}
  \correct the continuous functions
  \option the functions with $f(0) = 1$
  \correct the differentiable functions with $f' = f$
  \option the functions with $f(x) \ge 0$ for all $x$
\end{options}
\begin{solution}
Sums and multiples of continuous functions are continuous, and if $f' = f$ and $g' = g$, then $(af + bg)' = af + bg$. The functions with $f(0) = 1$ do not include 0, and the nonnegative functions are not closed under multiplication by $-1$.
\end{solution}
\end{question}

\begin{question}[id=vs-7]
If $W_1$ and $W_2$ are subspaces of a vector space $V$, which of the following is always a subspace?
\begin{options}
  \option $W_1 \cup W_2$
  \correct $W_1 \cap W_2$
  \option $V \setminus W_1$
  \option $\{w \in V : w \notin W_2\} \cup \{0\}$
\end{options}
\begin{solution}
If $u, v$ lie in both subspaces, so does $au + bv$. The union of two lines through 0 in $\mathbb{R}^2$ is not a subspace, and $V \setminus W_1$ does not contain 0.
\end{solution}
\end{question}

\begin{question}[id=vs-8, marks=2]
What is the dimension of the subspace of $M_3(\mathbb{R})$ consisting of the matrices $A$ with $A^T = A$ and $\operatorname{tr} A = 0$?
\answer{5}
\begin{solution}
The symmetric $3 \times 3$ matrices form a space of dimension 6 (choose the 3 diagonal entries and the 3 entries above the diagonal). The trace is a nonzero linear functional on it, so its kernel has dimension $6 - 1 = 5$.
\end{solution}
\end{question}

\begin{question}[id=vs-9, type=msq, marks=2]
Which of the following sets span $\mathbb{R}^3$?
\begin{options}
  \correct $\{(1, 0, 0),\, (0, 1, 0),\, (1, 1, 1)\}$
  \option $\{(1, 1, 0),\, (0, 1, 1),\, (1, 2, 1)\}$
  \correct $\{(1, 2, 3),\, (4, 5, 6),\, (7, 8, 9),\, (1, 0, 0)\}$
  \option $\{(1, 0, 0),\, (0, 1, 0)\}$
\end{options}
\begin{solution}
(A): the matrix with these columns is triangular with determinant 1. (B): $(1, 2, 1) = (1, 1, 0) + (0, 1, 1)$, so the span is only a plane. (C): $(1, 2, 3), (4, 5, 6), (1, 0, 0)$ already span $\mathbb{R}^3$, since $\det \begin{bmatrix} 1 & 4 & 1 \\ 2 & 5 & 0 \\ 3 & 6 & 0 \end{bmatrix} = -3 \ne 0$. (D): two vectors cannot span $\mathbb{R}^3$.
\end{solution}
\end{question}

% ==============================================================================================
\topic{basis}{Linear independence, basis and dimension}{Week 5}

\begin{question}[id=basis-1]
For which value of $k$ are the vectors $(1, 1, 0), (0, 1, 1), (1, 0, k)$ linearly dependent?
\begin{options}
  \option $k = 0$
  \option $k = 1$
  \correct $k = -1$
  \option for no value of $k$
\end{options}
\begin{solution}
They are dependent exactly when $\det \begin{bmatrix} 1 & 0 & 1 \\ 1 & 1 & 0 \\ 0 & 1 & k \end{bmatrix} = k + 1$ is 0, that is, $k = -1$. Indeed $(1,1,0) - (0,1,1) = (1, 0, -1)$.
\end{solution}
\end{question}

\begin{question}[id=basis-2, type=msq, marks=2]
Let $v_1, v_2, v_3$ be linearly independent vectors in a real vector space. Which of the following sets are also linearly independent?
\begin{options}
  \correct $\{v_1,\; v_1 + v_2,\; v_1 + v_2 + v_3\}$
  \option $\{v_1 - v_2,\; v_2 - v_3,\; v_3 - v_1\}$
  \correct $\{v_1 + v_2,\; v_2 + v_3,\; v_3 + v_1\}$
  \option $\{v_1,\; v_2,\; v_1 + v_2\}$
\end{options}
\begin{solution}
Write each set in terms of $v_1, v_2, v_3$; it is independent exactly when the coefficient matrix is invertible. (A): triangular with 1s on the diagonal, determinant 1. (B): the three vectors add up to 0. (C): $\det \begin{bmatrix} 1 & 0 & 1 \\ 1 & 1 & 0 \\ 0 & 1 & 1 \end{bmatrix} = 2 \ne 0$. (D): the third vector is the sum of the first two.
\end{solution}
\end{question}

\begin{question}[id=basis-3]
What is the dimension of $\operatorname{span}\{(1, 2, 1), (2, 1, 0), (3, 3, 1), (1, -1, -1)\}$ in $\mathbb{R}^3$?
\answer{2}
\begin{solution}
$(3, 3, 1) = (1, 2, 1) + (2, 1, 0)$ and $(1, -1, -1) = (2, 1, 0) - (1, 2, 1)$, while $(1, 2, 1)$ and $(2, 1, 0)$ are not multiples of each other. So the span has dimension 2.
\end{solution}
\end{question}

\begin{question}[id=basis-4]
A set of 5 vectors in $\mathbb{R}^4$
\begin{options}
  \option is always linearly independent
  \correct is always linearly dependent
  \option always spans $\mathbb{R}^4$
  \option never spans $\mathbb{R}^4$
\end{options}
\begin{solution}
An independent set in $\mathbb{R}^4$ has at most 4 vectors. Five vectors may or may not span $\mathbb{R}^4$ (the standard basis plus one more vector does).
\end{solution}
\end{question}

\begin{question}[id=basis-5, marks=2]
Which of the following is a basis of $\mathbb{P}_2$, the real polynomials of degree at most 2?
\begin{options}
  \correct $\{1,\; x,\; x^2 - x\}$
  \option $\{1 + x,\; 1 - x,\; 2\}$
  \option $\{x,\; x^2,\; x^3\}$
  \option $\{1,\; x^2\}$
\end{options}
\begin{solution}
(A): $x^2 = (x^2 - x) + x$, so the set spans $\mathbb{P}_2$, and 3 spanning vectors in a space of dimension 3 form a basis. (B) is dependent: $2 = (1 + x) + (1 - x)$. $x^3 \notin \mathbb{P}_2$, and 2 vectors cannot span a space of dimension 3.
\end{solution}
\end{question}

\begin{question}[id=basis-6, marks=2]
Let $W = \{(x_1, x_2, x_3, x_4) \in \mathbb{R}^4 : x_1 + x_2 = x_3 + x_4 \text{ and } x_1 = x_3\}$. Then $\dim W =$
\answer{2}
\begin{solution}
$W$ is the null space of $\begin{bmatrix} 1 & 1 & -1 & -1 \\ 1 & 0 & -1 & 0 \end{bmatrix}$, whose two rows are independent. So $\dim W = 4 - 2 = 2$; a basis is $(1, 0, 1, 0), (0, 1, 0, 1)$.
\end{solution}
\end{question}

\begin{question}[id=basis-7]
Let $V$ be a vector space of dimension $n \ge 2$. Which statement is \emph{false}?
\begin{options}
  \option Any $n$ linearly independent vectors form a basis of $V$.
  \option Any $n$ vectors that span $V$ form a basis of $V$.
  \option Any $n + 1$ vectors in $V$ are linearly dependent.
  \correct Any $n - 1$ vectors in $V$ are linearly independent.
\end{options}
\begin{solution}
(D) is false: for example, $n - 1$ copies of the same nonzero vector are dependent. (A), (B) and (C) are standard facts (the Basis Theorem).
\end{solution}
\end{question}

\begin{question}[id=basis-8, type=msq, marks=2]
Let $U$ and $W$ be subspaces of $\mathbb{R}^5$ with $\dim U = 3$ and $\dim W = 4$. Which values are possible for $\dim(U \cap W)$?
\begin{options}
  \option 1
  \correct 2
  \correct 3
  \option 4
\end{options}
\begin{solution}
$\dim(U \cap W) = \dim U + \dim W - \dim(U + W) \ge 3 + 4 - 5 = 2$, and $U \cap W \subseteq U$ gives $\dim(U \cap W) \le 3$. Both 2 and 3 occur (3 when $U \subseteq W$).
\end{solution}
\end{question}

\begin{question}[id=basis-9, marks=2]
Let $B = \begin{bmatrix} 1 & 1 \\ 0 & 1 \end{bmatrix}$. What is the dimension of the space of real $2 \times 2$ matrices $A$ with $AB = BA$?
\answer{2}
\begin{solution}
For $A = \begin{bmatrix} a & b \\ c & d \end{bmatrix}$, $AB = \begin{bmatrix} a & a + b \\ c & c + d \end{bmatrix}$ and $BA = \begin{bmatrix} a + c & b + d \\ c & d \end{bmatrix}$. These are equal exactly when $c = 0$ and $a = d$. So the space consists of the matrices $\begin{bmatrix} a & b \\ 0 & a \end{bmatrix}$, with basis $I$ and $\begin{bmatrix} 0 & 1 \\ 0 & 0 \end{bmatrix}$: its dimension is 2.
\end{solution}
\end{question}

% ==============================================================================================
\topic{rank}{Rank, nullity and the fundamental subspaces}{Week 6}

\begin{question}[id=rank-1]
The rank of $\begin{bmatrix} 1 & 2 & 3 \\ 4 & 5 & 6 \\ 7 & 8 & 9 \end{bmatrix}$ is
\answer{2}
\begin{solution}
Row 3 $-$ row 2 $=$ row 2 $-$ row 1 $= (3, 3, 3)$, so row 3 is $2 \cdot$ row 2 $-$ row 1, while rows 1 and 2 are independent. The rank is 2.
\end{solution}
\end{question}

\begin{question}[id=rank-2]
If $A$ is a $5 \times 7$ matrix of rank 4, then $\dim \operatorname{Nul} A$ is
\begin{options}
  \option 1
  \option 2
  \correct 3
  \option 4
\end{options}
\begin{solution}
Rank--nullity: $\operatorname{rank} A + \dim \operatorname{Nul} A = 7$, the number of columns. So the nullity is $7 - 4 = 3$.
\end{solution}
\end{question}

\begin{question}[id=rank-3, marks=2]
Let $A$ be an $m \times n$ matrix with $m \ne n$. Which of the following is always true?
\begin{options}
  \correct $\operatorname{rank} A = \operatorname{rank} A^T$
  \option $\operatorname{Nul} A = \operatorname{Nul} A^T$
  \option $\operatorname{Col} A = \operatorname{Row} A$
  \option $\operatorname{rank} A = \min(m, n)$
\end{options}
\begin{solution}
The row space and the column space always have the same dimension. (B) and (C) compare subspaces of $\mathbb{R}^n$ and $\mathbb{R}^m$, which are different spaces; (D) fails for the zero matrix.
\end{solution}
\end{question}

\begin{question}[id=rank-4, marks=2]
What is the dimension of the null space of
\[
  A = \begin{bmatrix} 1 & 2 & 0 & 1 & 3 \\ 2 & 4 & 1 & 3 & 7 \\ 1 & 2 & 1 & 2 & 4 \\ 0 & 0 & 1 & 1 & 1 \end{bmatrix}?
\]
\answer{3}
\begin{solution}
$R_2 - 2R_1$, $R_3 - R_1$ and $R_4$ are all equal to $(0, 0, 1, 1, 1)$, so the echelon form has 2 nonzero rows: rank 2. The nullity is $5 - 2 = 3$.
\end{solution}
\end{question}

\begin{question}[id=rank-5, type=msq, marks=2]
Let $A$ be a real $3 \times 3$ matrix of rank 2. Which of the following are true?
\begin{options}
  \correct $\det A = 0$
  \correct $Ax = 0$ has a nontrivial solution.
  \option $Ax = b$ is consistent for every $b \in \mathbb{R}^3$.
  \correct 0 is an eigenvalue of $A$.
\end{options}
\begin{solution}
Rank 2 < 3 means $A$ is not invertible: $\det A = 0$, $\operatorname{Nul} A \ne \{0\}$, and $Ax = 0\cdot x$ has nonzero solutions, so 0 is an eigenvalue. The column space is only a plane, so (C) fails.
\end{solution}
\end{question}

\begin{question}[id=rank-6]
What is the rank of the $4 \times 4$ matrix all of whose entries are 1?
\answer{1}
\begin{solution}
All four rows are equal and nonzero, so the row space is spanned by one vector: the rank is 1.
\end{solution}
\end{question}

\begin{question}[id=rank-7, marks=2]
Let $A$ be $4 \times 3$ and $B$ be $3 \times 4$. The $4 \times 4$ matrix $AB$
\begin{options}
  \option is always invertible
  \correct is never invertible
  \option is invertible if and only if $A$ and $B$ have rank 3
  \option has rank 4 whenever $A$ has rank 3
\end{options}
\begin{solution}
$\operatorname{Col}(AB) \subseteq \operatorname{Col} A$, so $\operatorname{rank}(AB) \le \operatorname{rank} A \le 3 < 4$. Hence $AB$ is never invertible.
\end{solution}
\end{question}

\begin{question}[id=rank-8, marks=2]
For which value of $k$ does $A = \begin{bmatrix} 1 & 1 & 1 \\ 1 & 2 & 3 \\ 1 & 3 & k \end{bmatrix}$ have rank 2?
\answer{5}
\begin{solution}
$\det A = 1(2k - 9) - 1(k - 3) + 1(3 - 2) = k - 5$. For $k = 5$ the determinant is 0, while the first two rows are independent, so the rank is 2; for $k \ne 5$ the rank is 3.
\end{solution}
\end{question}

\begin{question}[id=rank-9, marks=2]
Let $A$ be a real $5 \times 5$ matrix with $A^2 = 0$. The largest possible rank of $A$ is
\answer{2}
\begin{solution}
$A(Ax) = 0$ for every $x$, so $\operatorname{Col} A \subseteq \operatorname{Nul} A$ and $\operatorname{rank} A \le 5 - \operatorname{rank} A$, that is, $\operatorname{rank} A \le 2$. Rank 2 occurs: let $A$ send $e_3 \mapsto e_1$, $e_4 \mapsto e_2$ and the other standard basis vectors to 0.
\end{solution}
\end{question}

\begin{question}[id=rank-10, type=msq, marks=2]
Let $A$ and $B$ be real $3 \times 3$ matrices of rank 1. Which values are possible for $\operatorname{rank}(A + B)$?
\begin{options}
  \correct 0
  \correct 1
  \correct 2
  \option 3
\end{options}
\begin{solution}
$\operatorname{Col}(A + B) \subseteq \operatorname{Col} A + \operatorname{Col} B$, so $\operatorname{rank}(A + B) \le 1 + 1 = 2$. All of 0, 1 and 2 occur: take $B = -A$, then $B = A$, then $A = E_{11}$ and $B = E_{22}$ (the matrix units).
\end{solution}
\end{question}

% ==============================================================================================
\topic{lt}{Linear transformations and change of basis}{Weeks 7--8}

\begin{question}[id=lt-1, type=msq]
Which of the following maps $T : \mathbb{R}^2 \to \mathbb{R}^2$ are linear?
\begin{options}
  \correct $T(x, y) = (x + y,\; x - y)$
  \option $T(x, y) = (x + 1,\; y)$
  \option $T(x, y) = (xy,\; x)$
  \correct $T(x, y) = (2y,\; 0)$
\end{options}
\begin{solution}
(A) and (D) are given by matrices. (B) sends 0 to $(1, 0)$, so it is not linear. (C): $T(2, 2) = (4, 2)$ but $2T(1, 1) = (2, 2)$.
\end{solution}
\end{question}

\begin{question}[id=lt-2]
Let $T : \mathbb{R}^5 \to \mathbb{R}^3$ be a linear map whose range is all of $\mathbb{R}^3$. Then $\dim \ker T =$
\answer{2}
\begin{solution}
Rank--nullity: $\dim \ker T + \dim \operatorname{range} T = 5$, so $\dim \ker T = 5 - 3 = 2$.
\end{solution}
\end{question}

\begin{question}[id=lt-3]
The standard matrix of $T(x, y) = (2x + y,\; x - 3y)$ is
\begin{options}
  \correct $\begin{bmatrix} 2 & 1 \\ 1 & -3 \end{bmatrix}$
  \option $\begin{bmatrix} 2 & 1 \\ 1 & 3 \end{bmatrix}$
  \option $\begin{bmatrix} 1 & 2 \\ -3 & 1 \end{bmatrix}$
  \option $\begin{bmatrix} 2 & -3 \\ 1 & 1 \end{bmatrix}$
\end{options}
\begin{solution}
The columns are $T(e_1) = (2, 1)$ and $T(e_2) = (1, -3)$.
\end{solution}
\end{question}

\begin{question}[id=lt-4, marks=2]
Let $T : \mathbb{P}_2 \to \mathbb{P}_2$ be differentiation, $T(p) = p'$. Which is true?
\begin{options}
  \option $T$ is invertible.
  \correct $\ker T$ is the set of constant polynomials and $\operatorname{rank} T = 2$.
  \option $T$ is onto.
  \option $T^2 = 0$.
\end{options}
\begin{solution}
$p' = 0$ exactly for constants, so $\ker T$ has dimension 1 and $\operatorname{rank} T = 3 - 1 = 2$ (the range is $\mathbb{P}_1$). So $T$ is neither invertible nor onto, and $T^2(x^2) = 2 \ne 0$.
\end{solution}
\end{question}

\begin{question}[id=lt-5, marks=2]
What is the standard matrix of the map of $\mathbb{R}^2$ that rotates by $90^\circ$ anticlockwise and \emph{then} reflects in the $x$-axis?
\begin{options}
  \option $\begin{bmatrix} 0 & 1 \\ 1 & 0 \end{bmatrix}$
  \correct $\begin{bmatrix} 0 & -1 \\ -1 & 0 \end{bmatrix}$
  \option $\begin{bmatrix} 0 & 1 \\ -1 & 0 \end{bmatrix}$
  \option $\begin{bmatrix} 0 & -1 \\ 1 & 0 \end{bmatrix}$
\end{options}
\begin{solution}
It is $FR$ with $R = \begin{bmatrix} 0 & -1 \\ 1 & 0 \end{bmatrix}$ and $F = \begin{bmatrix} 1 & 0 \\ 0 & -1 \end{bmatrix}$: $FR = \begin{bmatrix} 0 & -1 \\ -1 & 0 \end{bmatrix}$. Check: $e_1 \mapsto (0, 1) \mapsto (0, -1)$. (It is the reflection in the line $y = -x$.)
\end{solution}
\end{question}

\begin{question}[id=lt-6]
Similar matrices $A$ and $B = P^{-1}AP$ always have the same
\begin{options}
  \option eigenvectors
  \correct characteristic polynomial
  \option reduced row echelon form
  \option diagonal entries
\end{options}
\begin{solution}
$\det(B - \lambda I) = \det(P^{-1}(A - \lambda I)P) = \det(A - \lambda I)$. The others can change. If $Av = \lambda v$, then $B(P^{-1}v) = \lambda P^{-1}v$, so the eigenvectors move. $\begin{bmatrix} 0 & 1 \\ 1 & 0 \end{bmatrix}$ and $\begin{bmatrix} 1 & 0 \\ 0 & -1 \end{bmatrix}$ are similar (both are diagonalizable with eigenvalues $\pm 1$) but have different diagonals. $\begin{bmatrix} 0 & 1 \\ 0 & 0 \end{bmatrix}$ and $\begin{bmatrix} 0 & 0 \\ 1 & 0 \end{bmatrix}$ are similar (swap the two basis vectors) but have different reduced echelon forms.
\end{solution}
\end{question}

\begin{question}[id=lt-7, marks=2]
Let $T : \mathbb{R}^3 \to \mathbb{R}^3$, $T(x, y, z) = (x + y,\; y + z,\; x + 2y + z)$. The dimension of the range of $T$ is
\answer{2}
\begin{solution}
The standard matrix has rows $(1,1,0)$, $(0,1,1)$, $(1,2,1)$. The third row is the sum of the first two, which are independent, so the rank, the dimension of the range, is 2.
\end{solution}
\end{question}

\begin{question}[id=lt-8, marks=2]
Let $\mathcal{B} = \{(1, 1), (1, -1)\}$, a basis of $\mathbb{R}^2$. The coordinate vector $[x]_{\mathcal{B}}$ of $x = (3, 1)$ is
\begin{options}
  \correct $(2, 1)$
  \option $(1, 2)$
  \option $(3, 1)$
  \option $(4, 2)$
\end{options}
\begin{solution}
Solve $a(1, 1) + b(1, -1) = (3, 1)$: $a + b = 3$, $a - b = 1$, so $a = 2$, $b = 1$.
\end{solution}
\end{question}

\begin{question}[id=lt-9, marks=2]
Let $T : \mathbb{R}^2 \to \mathbb{R}^2$ be the reflection in the line $y = x$. The matrix of $T$ relative to the basis $\mathcal{B} = \{(1, 1), (1, -1)\}$ is
\begin{options}
  \option $\begin{bmatrix} 0 & 1 \\ 1 & 0 \end{bmatrix}$
  \correct $\begin{bmatrix} 1 & 0 \\ 0 & -1 \end{bmatrix}$
  \option $\begin{bmatrix} 1 & 0 \\ 0 & 1 \end{bmatrix}$
  \option $\begin{bmatrix} -1 & 0 \\ 0 & 1 \end{bmatrix}$
\end{options}
\begin{solution}
$T(1, 1) = (1, 1)$ and $T(1, -1) = (-1, 1) = -(1, -1)$. The columns of $[T]_{\mathcal{B}}$ are the $\mathcal{B}$-coordinates of these images: $(1, 0)$ and $(0, -1)$. (Option A is the standard matrix of $T$; the two matrices are similar.)
\end{solution}
\end{question}

% ==============================================================================================
\topic{det}{Determinants}{Week 9}

\begin{question}[id=det-1]
$\det \begin{bmatrix} 2 & 0 & 0 \\ 5 & 3 & 0 \\ 7 & 1 & 4 \end{bmatrix} =$
\answer{24}
\begin{solution}
The matrix is lower triangular, so the determinant is the product of the diagonal entries: $2 \cdot 3 \cdot 4 = 24$.
\end{solution}
\end{question}

\begin{question}[id=det-2]
If $A$ is a $3 \times 3$ matrix with $\det A = 2$, then $\det(3A)$ is
\begin{options}
  \option 6
  \option 18
  \correct 54
  \option 8
\end{options}
\begin{solution}
Each of the 3 rows is multiplied by 3, so $\det(3A) = 3^3 \det A = 27 \cdot 2 = 54$.
\end{solution}
\end{question}

\begin{question}[id=det-3, marks=2]
If $A$ is a $4 \times 4$ matrix with $\det A = 3$, then $\det(2A^{-1}A^T) =$
\answer{16}
\begin{solution}
$\det(2A^{-1}A^T) = 2^4 \det(A^{-1}) \det(A^T) = 16 \cdot \tfrac{1}{3} \cdot 3 = 16$.
\end{solution}
\end{question}

\begin{question}[id=det-4, type=msq, marks=2]
For $n \times n$ matrices $A$ and $B$, which are always true?
\begin{options}
  \correct $\det(AB) = \det A \det B$
  \option $\det(A + B) = \det A + \det B$
  \correct $\det A^T = \det A$
  \option $\det(-A) = -\det A$
\end{options}
\begin{solution}
(A) and (C) are standard. (B) fails for $A = B = I_2$: $4 \ne 2$. (D): $\det(-A) = (-1)^n \det A$, which differs from $-\det A$ when $n$ is even and $\det A \ne 0$.
\end{solution}
\end{question}

\begin{question}[id=det-5, marks=2]
$\det \begin{bmatrix} 1 & 1 & 1 \\ 1 & 2 & 4 \\ 1 & 3 & 9 \end{bmatrix} =$
\answer{2}
\begin{solution}
A Vandermonde matrix in $1, 2, 3$: the determinant is $(2 - 1)(3 - 1)(3 - 2) = 2$. (Or row reduce: $R_2 - R_1 = (0, 1, 3)$, $R_3 - R_1 = (0, 2, 8)$, then $R_3 - 2R_2 = (0, 0, 2)$; the product of the pivots is 2.)
\end{solution}
\end{question}

\begin{question}[id=det-6]
Interchanging two rows of a square matrix
\begin{options}
  \option does not change its determinant
  \correct multiplies its determinant by $-1$
  \option multiplies its determinant by 2
  \option makes its determinant 0
\end{options}
\begin{solution}
A row interchange changes the sign of the determinant; a row replacement $R_i \to R_i + cR_j$ does not change it; scaling a row by $c$ multiplies it by $c$.
\end{solution}
\end{question}

\begin{question}[id=det-7, marks=2]
The area of the parallelogram with vertices $(0, 0), (3, 1), (1, 2), (4, 3)$ is
\begin{options}
  \correct 5
  \option 7
  \option 2.5
  \option 10
\end{options}
\begin{solution}
It is spanned by $(3, 1)$ and $(1, 2)$, so its area is $\left|\det \begin{bmatrix} 3 & 1 \\ 1 & 2 \end{bmatrix}\right| = |6 - 1| = 5$.
\end{solution}
\end{question}

\begin{question}[id=det-8, marks=2]
Find the value of $x \ne 1$ for which $\det \begin{bmatrix} x & 1 & 1 \\ 1 & x & 1 \\ 1 & 1 & x \end{bmatrix} = 0$.
\answer{-2}
\begin{solution}
The matrix is $(x - 1)I + J$, where $J$ has all entries 1 and eigenvalues $3, 0, 0$. So its eigenvalues are $x + 2, x - 1, x - 1$ and the determinant is $(x + 2)(x - 1)^2$. For $x \ne 1$ it vanishes at $x = -2$.
\end{solution}
\end{question}

\begin{question}[id=det-9]
Use Cramer's rule to solve $2x + y = 7,\; x + 3y = 11$. Then $y =$
\answer{3}
\begin{solution}
$y = \dfrac{\det \begin{bmatrix} 2 & 7 \\ 1 & 11 \end{bmatrix}}{\det \begin{bmatrix} 2 & 1 \\ 1 & 3 \end{bmatrix}} = \dfrac{22 - 7}{6 - 1} = 3$, and similarly $x = \dfrac{21 - 11}{5} = 2$.
\end{solution}
\end{question}

% ==============================================================================================
\topic{eig}{Eigenvalues, eigenvectors and diagonalization}{Weeks 10--11}

\begin{question}[id=eig-1]
The largest eigenvalue of $\begin{bmatrix} 2 & 1 \\ 1 & 2 \end{bmatrix}$ is
\answer{3}
\begin{solution}
$\det(A - \lambda I) = (2 - \lambda)^2 - 1 = (\lambda - 1)(\lambda - 3)$, so the eigenvalues are 1 and 3. (Eigenvectors $(1, -1)$ and $(1, 1)$.)
\end{solution}
\end{question}

\begin{question}[id=eig-2]
If $\lambda$ is an eigenvalue of an invertible matrix $A$, then an eigenvalue of $A^{-1}$ is
\begin{options}
  \option $\lambda$
  \option $-\lambda$
  \correct $1/\lambda$
  \option $\lambda^2$
\end{options}
\begin{solution}
From $Av = \lambda v$ (and $\lambda \ne 0$, as $A$ is invertible), multiply by $A^{-1}$: $v = \lambda A^{-1}v$, so $A^{-1}v = \tfrac{1}{\lambda} v$.
\end{solution}
\end{question}

\begin{question}[id=eig-3, marks=2]
Which of the following matrices is \emph{not} diagonalizable?
\begin{options}
  \correct $\begin{bmatrix} 1 & 1 \\ 0 & 1 \end{bmatrix}$
  \option $\begin{bmatrix} 1 & 0 \\ 0 & 1 \end{bmatrix}$
  \option $\begin{bmatrix} 2 & 1 \\ 1 & 2 \end{bmatrix}$
  \option $\begin{bmatrix} 1 & 2 \\ 0 & 3 \end{bmatrix}$
\end{options}
\begin{solution}
For (A), the only eigenvalue is 1 (with multiplicity 2), but $A - I = \begin{bmatrix} 0 & 1 \\ 0 & 0 \end{bmatrix}$ has rank 1, so the eigenspace is only 1-dimensional. (C) is symmetric and (D) has two distinct eigenvalues, 1 and 3.
\end{solution}
\end{question}

\begin{question}[id=eig-4, marks=2]
Let $A = \begin{bmatrix} 4 & -1 & 6 \\ 2 & 1 & 6 \\ 2 & -1 & 8 \end{bmatrix}$, which has eigenvalue 2. What is the dimension of the eigenspace of $A$ for $\lambda = 2$?
\answer{2}
\begin{solution}
$A - 2I = \begin{bmatrix} 2 & -1 & 6 \\ 2 & -1 & 6 \\ 2 & -1 & 6 \end{bmatrix}$ has rank 1, so its null space, the eigenspace, has dimension $3 - 1 = 2$. (A basis: $(1, 2, 0)$ and $(-3, 0, 1)$.)
\end{solution}
\end{question}

\begin{question}[id=eig-5, type=msq, marks=2]
Let $A$ be a real $3 \times 3$ matrix with eigenvalues 1, 2 and 3. Which of the following are true?
\begin{options}
  \correct $A$ is diagonalizable.
  \correct $\det A = 6$
  \correct $\operatorname{tr} A = 6$
  \option $A - I$ is invertible.
\end{options}
\begin{solution}
Three distinct eigenvalues give three independent eigenvectors, so $A$ is diagonalizable. The determinant is the product and the trace the sum of the eigenvalues: $6$ and $6$. Since 1 is an eigenvalue, $A - I$ is singular.
\end{solution}
\end{question}

\begin{question}[id=eig-6]
The eigenvalues of $\begin{bmatrix} 0 & -1 \\ 1 & 0 \end{bmatrix}$ are
\begin{options}
  \option $1, -1$
  \correct $i, -i$
  \option $0, 0$
  \option $1, 1$
\end{options}
\begin{solution}
$\det(A - \lambda I) = \lambda^2 + 1$, so $\lambda = \pm i$. (The matrix is a rotation by $90^\circ$, which fixes no real direction.)
\end{solution}
\end{question}

\begin{question}[id=eig-7, marks=2]
If a $3 \times 3$ matrix $A$ has eigenvalues $1, -1, 2$, then $\det(A^2 + I) =$
\answer{20}
\begin{solution}
If $Av = \lambda v$, then $(A^2 + I)v = (\lambda^2 + 1)v$. So the eigenvalues of $A^2 + I$ are $2, 2, 5$, and its determinant is $2 \cdot 2 \cdot 5 = 20$.
\end{solution}
\end{question}

\begin{question}[id=eig-8, marks=2]
Let $A$ be a real $n \times n$ matrix with $A^2 = I$ and $A \ne \pm I$. Then
\begin{options}
  \option 1 is the only eigenvalue of $A$
  \correct $A$ is diagonalizable, with eigenvalues 1 and $-1$
  \option $A$ is not diagonalizable
  \option $A = 0$
\end{options}
\begin{solution}
Every $v$ splits as $v = \tfrac{1}{2}(v + Av) + \tfrac{1}{2}(v - Av)$, with $A$ acting as $+1$ on the first part and $-1$ on the second. So $\mathbb{R}^n$ is a sum of the two eigenspaces and $A$ is diagonalizable; as $A \ne \pm I$, both eigenvalues occur.
\end{solution}
\end{question}

\begin{question}[id=eig-9, type=msq, marks=2]
Let $A$ be an $n \times n$ matrix. Which of the following are always true?
\begin{options}
  \correct Eigenvectors for distinct eigenvalues are linearly independent.
  \correct $A$ and $A^T$ have the same eigenvalues.
  \option $A$ and $A^T$ have the same eigenvectors.
  \correct If 0 is an eigenvalue of $A$, then $A$ is not invertible.
\end{options}
\begin{solution}
(B): $\det(A^T - \lambda I) = \det(A - \lambda I)^T = \det(A - \lambda I)$. (C) fails: $\begin{bmatrix} 1 & 1 \\ 0 & 2 \end{bmatrix}$ has the eigenvector $(1, 0)$, which is not an eigenvector of its transpose. (A) and (D) are standard.
\end{solution}
\end{question}

\begin{question}[id=eig-10]
The sum of the eigenvalues of $\begin{bmatrix} 1 & 2 & 3 \\ 4 & 5 & 6 \\ 7 & 8 & 9 \end{bmatrix}$ is
\answer{15}
\begin{solution}
The sum of the eigenvalues (with multiplicity) is the trace: $1 + 5 + 9 = 15$.
\end{solution}
\end{question}

\begin{question}[id=eig-11, marks=2]
For $A = \begin{bmatrix} 1 & 2 \\ 3 & 4 \end{bmatrix}$, the matrix $A^2 - 5A$ equals
\begin{options}
  \correct $2I$
  \option $-2I$
  \option $0$
  \option $5I$
\end{options}
\begin{solution}
The characteristic polynomial is $\lambda^2 - (\operatorname{tr} A)\lambda + \det A = \lambda^2 - 5\lambda - 2$. By the Cayley--Hamilton theorem, $A^2 - 5A - 2I = 0$, so $A^2 - 5A = 2I$. (Check: $A^2 = \begin{bmatrix} 7 & 10 \\ 15 & 22 \end{bmatrix}$.)
\end{solution}
\end{question}

\begin{question}[id=eig-12]
How many linearly independent eigenvectors does $\begin{bmatrix} 2 & 1 & 0 \\ 0 & 2 & 0 \\ 0 & 0 & 3 \end{bmatrix}$ have?
\answer{2}
\begin{solution}
The eigenvalues are 2 (twice) and 3. $A - 2I = \begin{bmatrix} 0 & 1 & 0 \\ 0 & 0 & 0 \\ 0 & 0 & 1 \end{bmatrix}$ has rank 2, so the eigenspace for 2 has dimension 1, and so has the eigenspace for 3. So at most 2 eigenvectors are independent, for example $e_1$ and $e_3$.
\end{solution}
\end{question}

\begin{question}[id=eig-13, type=msq, marks=2]
Let $J$ be the $4 \times 4$ matrix all of whose entries are 1. Which of the following are true?
\begin{options}
  \correct 4 is an eigenvalue of $J$.
  \correct The eigenspace of $J$ for the eigenvalue 0 has dimension 3.
  \correct $J$ is diagonalizable.
  \option $J^2 = J$
\end{options}
\begin{solution}
$J(1, 1, 1, 1) = 4(1, 1, 1, 1)$. $J$ has rank 1, so $\operatorname{Nul} J$, the eigenspace for 0, has dimension 3. That gives $1 + 3 = 4$ independent eigenvectors, so $J$ is diagonalizable (it is also symmetric). Each entry of $J^2$ is $1 + 1 + 1 + 1 = 4$, so $J^2 = 4J \ne J$.
\end{solution}
\end{question}

\begin{question}[id=eig-14, type=msq, marks=2]
Let $A$ be a real $n \times n$ matrix with $A^2 = A$. Which of the following are always true?
\begin{options}
  \option $I - A$ is invertible.
  \correct Every eigenvalue of $A$ is 0 or 1.
  \correct $A$ is diagonalizable.
  \correct $\operatorname{rank} A = \operatorname{tr} A$
\end{options}
\begin{solution}
If $Av = \lambda v$ with $v \ne 0$, then $\lambda v = A^2v = \lambda^2 v$, so $\lambda$ is 0 or 1. Every vector is $v = Av + (v - Av)$ with $A(Av) = Av$ and $A(v - Av) = 0$, so $\mathbb{R}^n$ is the sum of the eigenspaces for 1 and 0, and $A$ is diagonalizable. Then $A$ is similar to $\operatorname{diag}(1, \dots, 1, 0, \dots, 0)$, and its rank and trace both equal the number of 1s. (A) fails whenever $A \ne 0$: if $Av \ne 0$, then $(I - A)(Av) = Av - A^2v = 0$.
\end{solution}
\end{question}

% ==============================================================================================
\topic{orth}{Orthogonality, projections and least squares}{Week 12}

\begin{question}[id=orth-1]
The distance between $u = (1, 2, 3)$ and $v = (3, 4, 4)$ is
\answer{3}
\begin{solution}
$\|u - v\| = \|(-2, -2, -1)\| = \sqrt{4 + 4 + 1} = 3$.
\end{solution}
\end{question}

\begin{question}[id=orth-2]
The vectors $(1, 2, -1)$ and $(2, k, 4)$ are orthogonal when $k =$
\begin{options}
  \option 0
  \correct 1
  \option -1
  \option 2
\end{options}
\begin{solution}
Their dot product is $2 + 2k - 4 = 2k - 2$, which is 0 for $k = 1$.
\end{solution}
\end{question}

\begin{question}[id=orth-3, marks=2]
The orthogonal projection of $y = (7, 6)$ onto the line spanned by $u = (4, 2)$ is $(a, b)$. Then $a + b =$
\answer{12}
\begin{solution}
$\hat{y} = \dfrac{y \cdot u}{u \cdot u}\, u = \dfrac{28 + 12}{16 + 4}(4, 2) = 2(4, 2) = (8, 4)$, so $a + b = 12$.
\end{solution}
\end{question}

\begin{question}[id=orth-4, type=msq, marks=2]
Let $Q$ be a real $n \times n$ matrix whose columns are orthonormal. Which of the following are true?
\begin{options}
  \correct $Q^T Q = I$
  \correct $Q^{-1} = Q^T$
  \correct $\|Qx\| = \|x\|$ for every $x \in \mathbb{R}^n$
  \option $\det Q = 1$
\end{options}
\begin{solution}
The entries of $Q^TQ$ are the dot products of the columns, so $Q^TQ = I$; as $Q$ is square, $Q^{-1} = Q^T$. Then $\|Qx\|^2 = x^TQ^TQx = \|x\|^2$. But $\det Q = \pm 1$: a reflection has determinant $-1$.
\end{solution}
\end{question}

\begin{question}[id=orth-5, marks=2]
The Gram--Schmidt process applied to $x_1 = (1, 1, 0)$, $x_2 = (1, 0, 1)$ gives $v_1 = x_1$ and $v_2 =$
\begin{options}
  \correct $(\tfrac{1}{2}, -\tfrac{1}{2}, 1)$
  \option $(1, -1, 1)$
  \option $(0, -1, 1)$
  \option $(\tfrac{1}{2}, \tfrac{1}{2}, 1)$
\end{options}
\begin{solution}
$v_2 = x_2 - \dfrac{x_2 \cdot v_1}{v_1 \cdot v_1} v_1 = (1, 0, 1) - \tfrac{1}{2}(1, 1, 0) = (\tfrac{1}{2}, -\tfrac{1}{2}, 1)$. Check: $v_2 \cdot v_1 = 0$.
\end{solution}
\end{question}

\begin{question}[id=orth-6, marks=2]
The least-squares solution $\hat{x} = (\hat{x}_1, \hat{x}_2)$ of $Ax = b$, with $A = \begin{bmatrix} 1 & 0 \\ 1 & 1 \\ 1 & 2 \end{bmatrix}$ and $b = (6, 0, 0)$, has $\hat{x}_1 =$
\answer{5}
\begin{solution}
Solve the normal equations $A^TA\hat{x} = A^Tb$: $\begin{bmatrix} 3 & 3 \\ 3 & 5 \end{bmatrix}\hat{x} = \begin{bmatrix} 6 \\ 0 \end{bmatrix}$. Subtracting gives $2\hat{x}_2 = -6$, so $\hat{x}_2 = -3$ and $\hat{x}_1 = 5$: the best line through $(0, 6), (1, 0), (2, 0)$ is $y = 5 - 3t$.
\end{solution}
\end{question}

\begin{question}[id=orth-7]
If $W$ is a subspace of $\mathbb{R}^n$, then $\dim W + \dim W^\perp$ equals
\begin{options}
  \option 0
  \correct $n$
  \option $2n$
  \option $\dim W$
\end{options}
\begin{solution}
$\mathbb{R}^n = W \oplus W^\perp$, so the dimensions add up to $n$.
\end{solution}
\end{question}

\begin{question}[id=orth-8, marks=2]
For an $m \times n$ matrix $A$, the orthogonal complement of $\operatorname{Row} A$ is
\begin{options}
  \option $\operatorname{Col} A$
  \correct $\operatorname{Nul} A$
  \option $\operatorname{Nul} A^T$
  \option $\operatorname{Col} A^T$
\end{options}
\begin{solution}
$Ax = 0$ says exactly that $x$ is orthogonal to every row of $A$. So $(\operatorname{Row} A)^\perp = \operatorname{Nul} A$ (and $(\operatorname{Col} A)^\perp = \operatorname{Nul} A^T$).
\end{solution}
\end{question}

\begin{question}[id=orth-9, type=msq, marks=2]
Let $u \in \mathbb{R}^n$ ($n \ge 2$) be a unit vector and $P = I - 2uu^T$. Which of the following are true?
\begin{options}
  \correct $P$ is symmetric.
  \correct $P^2 = I$
  \option $\det P = 1$
  \correct $Pu = -u$
\end{options}
\begin{solution}
$P^T = I - 2(uu^T)^T = P$. Using $u^Tu = 1$: $Pu = u - 2u(u^Tu) = -u$, and $P^2 = I - 4uu^T + 4u(u^Tu)u^T = I$. $P$ fixes every vector orthogonal to $u$ (it is the reflection in the hyperplane $u^\perp$), so its eigenvalues are $-1$ once and $1$ with multiplicity $n - 1$, and $\det P = -1$.
\end{solution}
\end{question}

\begin{question}[id=orth-10]
Let $A = QR$, where the columns of $Q$ are orthonormal and $R$ is upper triangular. Then $R$ equals
\begin{options}
  \correct $Q^TA$
  \option $AQ^T$
  \option $QA$
  \option $A^TQ$
\end{options}
\begin{solution}
Orthonormal columns give $Q^TQ = I$, so $Q^TA = Q^TQR = R$. (By contrast, $A^TQ = R^TQ^TQ = R^T$, which is lower triangular.)
\end{solution}
\end{question}

% ==============================================================================================
\topic{sym}{Symmetric matrices, quadratic forms and the SVD}{Week 13}

\begin{question}[id=sym-1]
Every real symmetric matrix
\begin{options}
  \option is invertible
  \correct has only real eigenvalues
  \option has distinct eigenvalues
  \option is positive definite
\end{options}
\begin{solution}
The eigenvalues of a real symmetric matrix are real (spectral theorem). The zero matrix and $I$ show that the other statements can fail.
\end{solution}
\end{question}

\begin{question}[id=sym-2, marks=2]
The minimum of $Q(x, y) = 3x^2 + 2xy + 3y^2$ on the unit circle $x^2 + y^2 = 1$ is
\answer{2}
\begin{solution}
$Q(x) = x^TAx$ with $A = \begin{bmatrix} 3 & 1 \\ 1 & 3 \end{bmatrix}$, whose eigenvalues are 2 and 4. On unit vectors, $Q$ ranges over $[2, 4]$; the minimum 2 is attained at $\pm\tfrac{1}{\sqrt{2}}(1, -1)$.
\end{solution}
\end{question}

\begin{question}[id=sym-3, type=msq, marks=2]
Which of the following matrices are positive definite?
\begin{options}
  \correct $\begin{bmatrix} 2 & 1 \\ 1 & 2 \end{bmatrix}$
  \option $\begin{bmatrix} 1 & 2 \\ 2 & 1 \end{bmatrix}$
  \correct $\begin{bmatrix} 2 & 0 \\ 0 & 3 \end{bmatrix}$
  \option $\begin{bmatrix} 1 & 1 \\ 1 & 1 \end{bmatrix}$
\end{options}
\begin{solution}
A symmetric matrix is positive definite when all its eigenvalues are positive (for $2 \times 2$: $a > 0$ and $\det > 0$). Eigenvalues: (A) 1, 3; (B) $-1$, 3; (C) 2, 3; (D) 0, 2 (only semidefinite).
\end{solution}
\end{question}

\begin{question}[id=sym-4]
Eigenvectors of a real symmetric matrix that belong to distinct eigenvalues are
\begin{options}
  \option equal
  \correct orthogonal
  \option parallel
  \option linearly dependent
\end{options}
\begin{solution}
If $Au = \lambda u$, $Av = \mu v$ and $A^T = A$, then $\lambda\, u \cdot v = (Au) \cdot v = u \cdot (Av) = \mu\, u \cdot v$; with $\lambda \ne \mu$ this forces $u \cdot v = 0$.
\end{solution}
\end{question}

\begin{question}[id=sym-5, marks=2]
The largest singular value of $A = \begin{bmatrix} 3 & 0 \\ 4 & 5 \end{bmatrix}$ is (to two decimal places)
\answer{6.70 to 6.72}
\begin{solution}
$A^TA = \begin{bmatrix} 25 & 20 \\ 20 & 25 \end{bmatrix}$ has eigenvalues 45 and 5. The singular values are their square roots: $\sqrt{45} = 3\sqrt{5} \approx 6.71$ and $\sqrt{5}$.
\end{solution}
\end{question}

\begin{question}[id=sym-6, marks=2]
The quadratic form $Q(x) = x_1^2 + 4x_1x_2 + x_2^2$ is
\begin{options}
  \option positive definite
  \option negative definite
  \correct indefinite
  \option positive semidefinite
\end{options}
\begin{solution}
Its matrix is $\begin{bmatrix} 1 & 2 \\ 2 & 1 \end{bmatrix}$, with eigenvalues 3 and $-1$: one positive, one negative. Indeed $Q(1, 1) = 6 > 0$ and $Q(1, -1) = -2 < 0$.
\end{solution}
\end{question}

\begin{question}[id=sym-7]
For every real $m \times n$ matrix $A$, the matrix $A^TA$ is
\begin{options}
  \correct symmetric and positive semidefinite
  \option always invertible
  \option always positive definite
  \option skew-symmetric
\end{options}
\begin{solution}
$(A^TA)^T = A^TA$ and $x^TA^TAx = \|Ax\|^2 \ge 0$. It is positive definite only when the columns of $A$ are independent (the zero matrix shows it can be singular).
\end{solution}
\end{question}

\begin{question}[id=sym-8]
How many nonzero singular values does $A = \begin{bmatrix} 1 & 2 \\ 2 & 4 \\ 3 & 6 \end{bmatrix}$ have?
\answer{1}
\begin{solution}
The number of nonzero singular values is the rank. The second column is twice the first, so the rank is 1.
\end{solution}
\end{question}

\begin{question}[id=sym-9, marks=2]
What is the largest integer $k$ for which $\begin{bmatrix} 2 & k \\ k & 8 \end{bmatrix}$ is positive definite?
\answer{3}
\begin{solution}
A symmetric $2 \times 2$ matrix is positive definite exactly when its $(1,1)$ entry and its determinant are positive. Here $2 > 0$, and $16 - k^2 > 0$ exactly when $|k| < 4$. So the largest integer is $k = 3$.
\end{solution}
\end{question}

\end{document}
